(i) By Fact \claimref{standard-model.imported-clebsch-structure-seesaw}, under $g_{\lambda}$ the $126$ terms of the four Dirac mass matrices are $(w_{u}/\lambda)(\lambda f) = w_{u}f$, $-3(w_{u}/\lambda)(\lambda f) = -3w_{u}f$, and likewise with $w_{d}$; the $10$ and $120$ terms are untouched, so $M_{u}, M_{\nu}, M_{d}, M_{e}$ are invariant. $M_{R} = v_{R}f \mapsto \lambda v_{R}f$, and $m_{\nu} = -M_{\nu}(\lambda M_{R})^{-1}M_{\nu}^{\mathsf{T}} = \lambda^{-1}m_{\nu}$. Then $m_{\nu}m_{\nu}^{\dagger} \mapsto |\lambda|^{-2}m_{\nu}m_{\nu}^{\dagger}$, a positive multiple, which has the same eigenvectors and the same eigenvalue ratios; the Takagi values of $M_{R}$ (heavy masses) scale by $|\lambda|$ and those of $m_{\nu}$ by $|\lambda|^{-1}$. The light-neutrino unitary is the product of the charged-lepton left unitary (from $M_{e}M_{e}^{\dagger}$, invariant) and the eigenvectors of $m_{\nu}m_{\nu}^{\dagger}$ (invariant up to the phases that low-energy data do not fix), so the lepton mixing angles are invariant.

(ii) (EW) states $(M_{\nu})_{33} = 0$ and $(M_{u})_{33} = m_{t}$ in its family basis (equivalently $4w_{u}f_{33} = m_{t}$, since $h$ has no diagonal); both entries belong to invariant matrices, and $w_{u}f_{33}$ is invariant. Theorem \claimref{synthesis.hypothesis-u-under-ew-is-a-charged-fermion-condition}(i) reads Hypothesis (U), as its formal text states it, as a condition on the leading singular value of $M_{\nu}$, which is invariant. The chain of Proposition \claimref{standard-model.chain} solves for $(\ln M_{I}, \ln M_{U}, \alpha_{U}^{-1})$ from the three gauge couplings at $M_{Z}$ with the one-loop coefficient sums and the two-loop gauge matrix of its clause (1); no Yukawa coupling enters, so $M_{I}$, $g_{R}(M_{I})$ and $v_{R} = M_{I}/g_{R}$ (Proposition \claimref{standard-model.what-angle-asks-neutrino-sector}) are invariant. The eighteen observables of Proposition \claimref{standard-model.fits-record-level}: the nine charged-fermion masses and the four CKM numbers are functions of $M_{u}M_{u}^{\dagger}$, $M_{d}M_{d}^{\dagger}$ and $M_{e}M_{e}^{\dagger}$; the three lepton angles and $m_{3}/m_{2}$ are invariant by (i); $m_{3}$ scales by $|\lambda|^{-1}$. The rank-one seesaw of Proposition \claimref{neutrino.rank-one-seesaw-ordering-sum-double}, $m_{3} = m_{D}^{2}/M$, is the same scaling read along the heavy direction.

(iii) Let $D$ be data satisfying every stored equality and reproducing the record. Then $g_{\lambda}D$ satisfies every stored equality and reproduces every record value except $m_{3}$, which it reproduces exactly at $|\lambda| = 1$. So the stored equalities, with all record values except $m_{3}$, leave $|\lambda|$ free; the observed $m_{3}$ fixes it. With Theorem \claimref{synthesis.majorana-scale-is-126-doublet-vev}(ii), $(M_{R})_{33} = v_{R}f_{33} = m_{t}v_{R}/(4w_{u})$ moves as $\lambda$ along the orbit. The bounds: $|f_{33}| \le 4\pi$ gives $|(M_{R})_{33}| \le 4\pi v_{R}$; $|w_{u}| \le v$ gives $|(M_{R})_{33}| \ge m_{t}v_{R}/(4v)$. Their ratio is $16\pi v/m_{t} > 48 \cdot 174/100 > 83$ for $m_{t}(M_{I}) \le 100$ GeV (using $\pi > 3$); at $v_{R} = 4.2 \times 10^{10}$ GeV and $m_{t}(M_{I}) = 91$ GeV the range is $5.49 \times 10^{9}$ to below $5.28 \times 10^{11}$ GeV (using $\pi < 22/7$), the values of the stored proof. The odd towers' gravitational-seeding window under Premise (R) needs $M_{3} > 8.3 \times 10^{11}$ GeV at the largest content the seeding result considers (its restated value, with the derived production coefficients), above the range; that stored no-go rests on Hypothesis (U), which contradicts (EW), so it is quoted here as a number and not used. The normalisation of $w_{u}$ against the canonically normalised doublet, needed for $|w_{u}| \le v$, is cited from memory in that stored proof and is of order one.

(iv) The chain is a square system: three couplings at $M_{Z}$ determine its three unknowns. An independent determination of $v_{R}$ (hence of $M_{I} = g_{R}v_{R}$ at the chain's $g_{R}$) adds one equation and makes one of the three inputs, conventionally $\sin^{2}\theta_{W}$, a prediction; by (ii) $v_{R}$ is invariant under $g_{\lambda}$, so this determination says nothing about $\lambda$, and $M_{R} = \lambda$-dependent stays free. Conversely a determination of $M_{R}$ fixes $\lambda$ at given $v_{R}$ and adds no equation to the chain. A relation fixing $f$ (or $w_{u}$) against the gauge data would break $g_{\lambda}$ and would make both statements predictions; none is stored.

Electroweak normalisation. The doublet vevs satisfy $\sum_{a}(|v_{u}^{a}|^{2} + |v_{d}^{a}|^{2}) + |w_{u}|^{2} + |w_{d}|^{2} + (\text{$120$ terms}) = v^{2}$ in the canonical normalisation (from memory; not a stored relation), which $g_{\lambda}$ changes. The rescaling $(Y_{a}, v^{a}_{u}, v^{a}_{d}) \mapsto (\mu_{a}Y_{a}, v^{a}_{u}/\mu_{a}, v^{a}_{d}/\mu_{a})$, and $(h, c) \mapsto (\nu h, c/\nu)$, fixes every Dirac matrix and every relation of (i) to (iii), so composing it with $g_{\lambda}$ restores the sum. Under (EW) the $10$s carry about $\tfrac34 m_{t}$ in the $33$ entry, so perturbativity of $Y_{a}$ needs only $|v^{a}| \gtrsim 3m_{t}/(16\pi)$, a few GeV, and the interval of (iii) is essentially unchanged.

(v) The relations of (i) to (iii) are the Clebsch structure, (EW), the Dirac reading of (U), the chain, the fits and the stated inequalities; none contains $H$, $\Lambda$, $N$ or the rung unit $\mu$ of Premise \claimref{synthesis.rung-unit-second-free-number}. So $g_{\lambda}$, composed with the rescaling above when the electroweak normalisation is imposed, acts at fixed $N$ and $v/H$, and $M_{R}/M_{P}$ varies along an orbit on which both are constant. This shows that no stored relation ties $M_{R}$ to $N$ or $v/H$; the fits already have more couplings than observables (Proposition \claimref{standard-model.fits-record-level}), so on this route the Majorana scale is one more fitted input among several.

Scope. The symmetry is a symmetry of the stored relations. The two-loop gauge running with Yukawa contributions, not included in the store's chain, contains $\mathrm{Tr}(ff^{\dagger})$ and breaks $g_{\lambda}$; near the perturbative edge $|f| \sim 4\pi$ this is not small, so a two-loop chain with Yukawa terms would tie $M_{I}$ weakly to $f$. That over-determination is not computed here. The scalar potential, which would fix $w_{u}$ dynamically, is not in the store.

The check verifies (i) on complex rational instances with two $10$s, a symmetric $f$, an antisymmetric $h$ and generic complex vevs, for $\lambda \in \{7/3,\ 2 + I_{4},\ -1/5 + 3I_{4}/4\}$ (numbers, and $\C$, are read in $\mathrm{span}\{1, I_{4}\}$), including $\lambda M_{R}$, $\lambda^{-1}m_{\nu}$, the Gram matrix of $m_{\nu}$ up to $|\lambda|^{-2}$, and the scale-free invariants $\mathrm{tr}(H^{2})/\mathrm{tr}(H)^{2}$ and $\det H/\mathrm{tr}(H)^{3}$; the control that $f \mapsto \lambda f$ alone moves the Dirac mass matrices; (ii) with (EW) imposed along a heavy direction, the invariance of (EW), of the Dirac Gram matrix and the scaling of the heavy entry; the count seventeen invariant and one moving; and the interval of (iii).
